\section*{Chapitre 13 - Triangles}
\addcontentsline{toc}{section}{Chapitre 13 - Triangles}

\subsection*{I. Généralités sur les triangles}
\addcontentsline{toc}{subsection}{I. Généralités sur les triangles}

  \vspace{0.3cm}
\begin{Définition}
Un \textbf{triangle} est un polygone à trois côtés qui possède donc trois angles.
\end{Définition}
   \vspace{0.3cm}
\begin{Propriété}[c]\textbf{Inégalité triangulaire.}\\
Dans un triangle $ABC$ non plat, on a les inégalités suivantes :
\vspace{-0.5cm}
\begin{multicols}{3}
\begin{enumerate}
\item $AB < AC + CB$
\item $AC < AB + BC$
\item $BC < BA + AC$
\end{enumerate}
\end{multicols}
\end{Propriété}
 
  \vspace{0.3cm}
\begin{Propriété}[c]
Un triangle est constructible lorsque la longueur de son plus grand côté est strictement inférieure à la somme des longueurs des deux autres côtés.
\end{Propriété}
  \vspace{0.3cm}
 
\begin{Méthode} {\footnotesize Tracer un triangle connaissant les mesures des trois côtés. (LLS.fr/M6Methode28)} \\
Tracer le triangle MAT tel que MA = 5 cm, MT = 4 cm et AT = 6cm.

\vspace{6cm}

\end{Méthode}
  \vspace{0.3cm}
\begin{Méthode} {\footnotesize Tracer un triangle connaissant les mesures de deux côtés et d'un angle. (LLS.fr/M6Methode29) }\\
Tracer le triangle RST tel que RT = 5 cm, ST = 3 cm et $\widehat{RTS}$ = 60°.

\vspace{5cm}

\end{Méthode}

\newpage

\begin{Méthode} {\footnotesize Tracer un triangle connaissant les mesures de deux angles et d'un côté.}\\
Tracer le triangle EFG tel que EF = 7 cm, $\widehat{FEG}$ = 35° et $\widehat{EFG}$ = 40°.

\vspace{4.5cm}

\end{Méthode}

\begin{Propriété}[c]
Dans un triangle, la somme des angles est égale à $180$°.\\
Les trois angles d'un triangle sont donc supplémentaires.
\end{Propriété}
 

\begin{Exemple}
\begin{minipage}[c]{0.75\textwidth}
Calculer la mesure de $\widehat{ACB}$. 

\ligne

\ligne

\ligne

\end{minipage}\hfill  \begin{minipage}[c]{0.25\textwidth}
\begin{center}
\begin{tikzpicture}[scale=0.6]

% Définition des coordonnées des points
\coordinate (A) at (-2,3);
\coordinate (B) at (-4,0);
\coordinate (C) at (1,0);


% Tracé et coloriage des angles
% Angle en A (entre AB et AC)
% Direction AB depuis A: vers B-A = (-2,-3) 
% Direction AC depuis A: vers C-A = (3,-3)
\fill[Couleur8!30] (A) -- +(-0.5,-0.75) arc[start angle=236, end angle=315, radius=0.9] -- cycle;
\draw[Couleur8, thick] (A) +(-0.5,-0.75) arc[start angle=236, end angle=315, radius=0.9];
\node[Couleur8] at (-1.9,1.6) {$78$°};

% Angle en B (entre BA et BC)
% Vecteur BA: (2,3), angle ≈ 56°
% Vecteur BC: (5,0), angle = 0°
\fill[Couleur8!30] (B) -- +(0.8,0) arc[start angle=0, end angle=56, radius=0.8] -- cycle;
\draw[Couleur8, thick] (B) +(0.8,0) arc[start angle=0, end angle=56, radius=0.8];
\node[Couleur8] at (-2.7,0.5) {$60$°};

% Angle en C (entre CB et CA)
% Vecteur CB: (-5,0), angle = 180°
% Vecteur CA: (-3,3), angle ≈ 135°
\fill[Couleur8!30] (C) -- +(-0.8,0) arc[start angle=180, end angle=135, radius=0.8] -- cycle;
\draw[Couleur8, thick] (C) +(-0.8,0) arc[start angle=180, end angle=135, radius=0.8];
\node[Couleur8] at (-0.1,0.5) {?};

% Lignes du triangle
\draw[ultra thick, Couleur8] (A) -- (B);
\draw[ultra thick, Couleur8] (B) -- (C);
\draw[ultra thick, Couleur8] (C) -- (A);


% Étiquettes des points
\node[] at (A) {+};
\node[] at (B) {+};
\node[] at (C) {+};
\node[above] at (A) {A};
\node[left] at (B) {B};
\node[right] at (C) {C};

\end{tikzpicture}
\end{center}
\end{minipage}

\end{Exemple}

\begin{minipage}[c]{0.65\textwidth}
\begin{Définition}
Le \textbf{cercle circonscrit} à un triangle est le cercle passant par les trois sommets du triangle.
\end{Définition}
 
\begin{Propriété}[c]
Le centre du cercle circonscrit à un triangle est le point de concours des trois médiatrices de ce triangle.
\end{Propriété}
\end{minipage}\hfill  \begin{minipage}[c]{0.35\textwidth}
\begin{center}
\begin{tikzpicture}[scale=0.6]

% Points du triangle
\coordinate (A) at (-2,0);
\coordinate (B) at (2,1);
\coordinate (C) at (0,4);
\coordinate (D) at (-0.29,1.64);
\coordinate (E) at (-1,2);
\coordinate (F) at (0,0.5);
\coordinate (G) at (1,2.5);
\coordinate (I) at (2.98,3.82);
\coordinate (J) at (-3.55,-0.53);
\coordinate (K) at (0.67,-2.16);
\coordinate (L) at (-1.24,5.45);
\coordinate (M) at (-3.79,3.4);
\coordinate (N) at (3.22,-0.11);


% Centre du cercle (à la main pour simplifier)


% Triangle
\draw[thick, Couleur7] (A) -- (B) -- (C) -- cycle;


% Cercle circonscrit
\draw[Couleur6, very thick] (-0.29,1.64) circle (2.37);
\draw[Couleur6, very thick] (N) -- (M);
\draw[Couleur6, very thick] (K) -- (L);
\draw[Couleur6, very thick] (I) -- (J);

\node[Couleur7] at (A) {+};
\node[Couleur7] at (B) {+};
\node[Couleur7] at (C) {+};
\node[Couleur7] at (D) {+};

\node[right, Couleur6] at (L) {$(d_1)$};
\node[above right, Couleur6] at (M) {$(d_2)$};
\node[above, Couleur6] at (-3.55,-0.3) {$(d_3)$};

\node[below left,Couleur7] at (A) {A};
\node[below right,Couleur7] at (B) {B};
\node[above,Couleur7] at (C) {C};
\node[above,Couleur7] at (-0.1,1.7) {D};

% Pour AQ = QB (un trait perpendiculaire sur chaque segment)
\coordinate (MidAE) at ($(A)!0.5!(E)$);
\coordinate (MidEC) at ($(E)!0.5!(C)$);
\draw[thick, gray] ($(MidAE)!0.15cm!90:(E)$) -- ($(MidAE)!0.15cm!-90:(E)$);
\draw[thick, gray] ($(MidEC)!0.15cm!90:(C)$) -- ($(MidEC)!0.15cm!-90:(C)$);

% Pour AP = PB (un cercle sur chaque segment)  
\coordinate (MidCG) at ($(C)!0.5!(G)$);
\coordinate (MidGB) at ($(G)!0.5!(B)$);
\draw[thick, gray] (MidCG) circle (0.15cm);
\draw[thick, gray] (MidGB) circle (0.15cm);

% Pour AN = NB (une croix sur chaque segment)
\coordinate (MidAF) at ($(A)!0.5!(F)$);
\coordinate (MidFB) at ($(F)!0.5!(B)$);
\draw[thick, gray] ($(MidAF) + (-0.15,0)$) -- ($(MidAF) + (0.15,0)$);
\draw[thick, gray] ($(MidAF) + (0,-0.15)$) -- ($(MidAF) + (0,0.15)$);
\draw[thick, gray] ($(MidFB) + (-0.15,0)$) -- ($(MidFB) + (0.15,0)$);
\draw[thick, gray] ($(MidFB) + (0,-0.15)$) -- ($(MidFB) + (0,0.15)$);

% Petits carrés pour marquer les angles droits

% Angle droit en E (angle DEC)
\coordinate (E1) at ($(E)!0.3cm!(D)$);
\coordinate (E2) at ($(E)!0.3cm!(C)$);
\coordinate (E3) at ($(E1)+(E2)-(E)$);
\draw[thick, gray] (E1) -- (E3) -- (E2);

% Angle droit en G (angle DGB)
\coordinate (G1) at ($(G)!0.3cm!(D)$);
\coordinate (G2) at ($(G)!0.3cm!(B)$);
\coordinate (G3) at ($(G1)+(G2)-(G)$);
\draw[thick, gray] (G1) -- (G3) -- (G2);

% Angle droit en F (angle DFA)
\coordinate (F1) at ($(F)!0.3cm!(D)$);
\coordinate (F2) at ($(F)!0.3cm!(A)$);
\coordinate (F3) at ($(F1)+(F2)-(F)$);
\draw[thick, gray] (F1) -- (F3) -- (F2);


\end{tikzpicture}
\end{center}
\end{minipage}

 \vspace{-0.5cm}
\begin{Remarque}
\vspace{0.2cm}
Le centre du cercle circonscrit peut se trouver à l'extérieur du triangle.
\end{Remarque}

\begin{Exemple} 
Tracer le cercle circonscrit au triangle RST.
\begin{center}
\vspace{5cm}
\end{center}
\end{Exemple}

\newpage
\subsection*{II. Triangles rectangles}
\addcontentsline{toc}{subsection}{II. Triangles rectangles}

 
\begin{Définition}[c]
Un \textbf{triangle rectangle} est un triangle possédant un angle droit.\\
Le côté opposé à l'angle droit s'appelle l'\textbf{hypoténuse}.
\end{Définition}
 


\begin{Propriété}[c]
Dans un triangle rectangle, l'hypoténuse est toujours le plus grand côté.
\end{Propriété}

 

\begin{Exemple}
\begin{minipage}[c]{0.6\textwidth}
$DEF$ est un triangle rectangle en $F$. L'angle $\widehat{DFE}$ est un angle $\ldots\ldots$ , l'hypoténuse est le segment $\ldots\ldots$.

\ligne

\ligne

\ligne

\end{minipage}\hfill  \begin{minipage}[c]{0.35\textwidth}
\begin{center}
\begin{tikzpicture}[scale=1]

    
    % Côtés du triangle
    \draw[thick, Couleur8] (0,0) -- (3,0);  % Base DF
    \draw[thick, Couleur8] (3,0) -- (3,2);  % Hauteur FE
    \draw[thick, Couleur7, line width=1.5pt] (0,0) -- (3,2);  % Hypoténuse DE
    
    % Angle droit
    \draw[Couleur8] (2.7,0) -- (2.7,0.3) -- (3,0.3);
    
    % Arc pour l'angle de 30°
    \draw[Couleur8] (0.5,0) arc (0:33.69:0.5);
    
    % Labels des points
    \node[Couleur8] at (-0.2,-0.2) {D};
    \node[Couleur8] at (3.2,-0.2) {F};
    \node[Couleur8] at (3.3,2) {E};
    
    % Angle de 30°
    \node[Couleur8] at (1.1,0.3) {30°};
    
   % Arc pour l'angle de 60° en E
    \draw[Couleur8] (3,1.5) arc (270:213.43:0.5);
    
    % Point d'interrogation
    \node[Couleur8] at (2.6,1.2) {?};
    
    % Label "Hypoténuse" avec flèche
    \node[Couleur7] at (0.7,2) {Hypoténuse};
    \draw[->, Couleur7, thick] (0.7,1.7) -- (1.5,1.2);
    
\end{tikzpicture}
\end{center}
\end{minipage}

\end{Exemple}
 
 
 
\subsection*{III. Triangles isocèles et équilatéraux}
\addcontentsline{toc}{subsection}{III. Triangles isocèles et équilatéraux}

 
\begin{Définitions}[c]
Un triangle \textbf{isocèle} est un triangle possédant deux côtés de même longueur. Le sommet commun aux deux côtés de même longueur s'appelle le \textbf{sommet principal}. Le côté opposé au sommet principal s'appelle la \textbf{base}.
\end{Définitions}
 
\begin{Propriété}[c]
Dans un triangle isocèle, les angles à la base ont la même mesure.
\end{Propriété}
 
\begin{Exemple}
\begin{minipage}[c]{0.6\textwidth}
$ABC$ est isocèle en $\ldots$ (le sommet principal), sa base est le côté $\ldots$.\\
Les angles $\widehat{CAB}$ et $\widehat{ABC}$ ont \ligne

\end{minipage}\hfill  \begin{minipage}[c]{0.35\textwidth}

\begin{center}

\begin{tikzpicture}[scale=1]
	\coordinate (A) at (0,0);
	\coordinate (B) at (4,0);
	\coordinate (C) at (2,1.5);
    
    % Côtés du triangle
    \draw[ultra thick, Couleur7] (A) -- (B);  % Base AB
    \draw[very thick, Couleur8] (A) -- (C);  % Côté AC
    \draw[very thick, Couleur8] (B) -- (C);  % Côté BC
    

    
    % Labels des points
    \node[Couleur8,left] at (A) {A};
    \node[Couleur8,right] at (B) {B};
    \node[Couleur11,above] at (C) {C};
    
    % Label "Base"
    \node[Couleur7] at (2,-0.5) {Base};
    
    % Label "Sommet principal" avec flèche
    \node[Couleur11,left] at (0.5,1.75) {Sommet};
    \node[Couleur11,left] at (0.5,1.25) {principal};
    \draw[->, Couleur11, thick] (0.6,1.5) -- (1.8,1.5);
    
% Pour AC = CB (un trait perpendiculaire sur chaque segment)
\coordinate (MidAC) at ($(A)!0.5!(C)$);
\coordinate (MidCB) at ($(C)!0.5!(B)$);
\draw[thick, gray] ($(MidAC)!0.15cm!90:(C)$) -- ($(MidAC)!0.15cm!-90:(C)$);
\draw[thick, gray] ($(MidCB)!0.15cm!90:(B)$) -- ($(MidCB)!0.15cm!-90:(B)$);
    
\end{tikzpicture}
\end{center}
\end{minipage}

\end{Exemple}

\begin{Remarque}
Il est possible que l'angle situé au niveau du sommet principal d'un triangle isocèle soit un angle droit. Dans ce cas, on dit que c'est un triangle \textbf{isocèle rectangle}.
\end{Remarque}
 
\begin{minipage}[c]{0.70\textwidth}
\begin{Définition}[c]
Un triangle \textbf{équilatéral} est un triangle possédant trois côtés de même longueur.
\end{Définition}
 
\begin{Propriété}[c]
Dans tout triangle équilatéral, chaque angle mesure $60$°.
\end{Propriété}
\end{minipage}\hfill  \begin{minipage}[c]{0.3\textwidth}
\begin{center}

\begin{tikzpicture}[scale=0.8]
	\coordinate (A) at (0,0);
	\coordinate (B) at (4,0);
	\coordinate (C) at (2,3.5);
    
     % Angles du triangle
    % Angle en A
    \draw[thick, Couleur6!30] (A) ++(0.8,0) arc (0:60:0.8);

    
    % Angle en B
    \draw[thick, Couleur6!30] (B) ++(180:0.8) arc (180:120:0.8);

    
    % Angle en C
    \draw[thick, Couleur6!30] (C) ++(240:0.8) arc (240:300:0.8);

    
    
    % Côtés du triangle
    \draw[very thick, Couleur6] (A) -- (B);  % Base AB
    \draw[very thick, Couleur6] (A) -- (C);  % Côté AC
    \draw[very thick, Couleur6] (B) -- (C);  % Côté BC
    

    
    % Labels des points
    \node[Couleur6,left] at (A) {A};
    \node[Couleur6,right] at (B) {B};
    \node[Couleur6,above] at (C) {C};
    

    
    
% Pour AC = CB (un trait perpendiculaire sur chaque segment)
\coordinate (MidAC) at ($(A)!0.5!(C)$);
\coordinate (MidCB) at ($(C)!0.5!(B)$);
\coordinate (MidAB) at ($(A)!0.5!(B)$);
\draw[thick, gray] ($(MidAC)!0.15cm!90:(C)$) -- ($(MidAC)!0.15cm!-90:(C)$);
\draw[thick, gray] ($(MidCB)!0.15cm!90:(B)$) -- ($(MidCB)!0.15cm!-90:(B)$);
\draw[thick, gray] ($(MidAB)!0.15cm!90:(B)$) -- ($(MidAB)!0.15cm!-90:(B)$);


\end{tikzpicture}
\end{center}
\end{minipage}

